\documentclass[12pt]{article}
\usepackage[utf8]{inputenc}
\usepackage[T1]{fontenc}
\usepackage[margin=1in]{geometry}
\usepackage{setspace}
\usepackage{xcolor}
\usepackage[colorlinks=true, urlcolor=blue!60!black, linkcolor=blue!60!black]{hyperref}
\setlength{\parindent}{2em}
\setlength{\parskip}{1.25em}
\renewcommand{\baselinestretch}{1.0}
\begin{document}
\begin{center}
{\LARGE \textit{Murphy v.\ National Collegiate Athletic Association}, 584 U.S.\ 453 (2018)\par}
\vspace{0.8em}
{\large Case Brief --- Four-Page Edition\par}
\end{center}
\vspace{1em}

\section{Memory Jogger}

The anti-commandeering rule bars Congress from issuing direct orders to state legislatures---whether commanding them to enact laws or forbidding them to repeal or modify existing ones---so the federal Professional and Amateur Sports Protection Act (PASPA) fell in its entirety (\href{https://supreme.justia.com/cases/federal/us/584/16-476/}{Justia}).

\section{Detailed Case Facts}

PASPA (1992) made it unlawful for a State to ``sponsor, operate, advertise, promote, license, or authorize by law or compact'' sports-gambling schemes (\S~3702(1)), and barred private parties from operating such schemes pursuant to state law (\S~3702(2)). It did not make sports gambling a federal crime; instead, the Attorney General and sports organizations could sue for injunctions. Nevada, Delaware, Oregon, and Montana were grandfathered, and New Jersey had a one-year window it did not use. Facing casino competition, New Jersey voters in 2011 amended the state constitution, and the legislature legalized sports betting in 2012. After that law was enjoined, a 2014 Act repealed state prohibitions on sports wagering for those 21 and older at Atlantic City casinos and racetracks, leaving bans on wagers involving New Jersey college teams.

\section{Procedural History}

The NCAA and the NBA, NFL, NHL, and MLB sued Governor Christie. In \textit{Christie I}, the courts enjoined the 2012 Act and the Third Circuit rejected the anti-commandeering defense; certiorari was denied in 2014. In \textit{Christie II}, joined on the defense side by the New Jersey Thoroughbred Horsemen's Association, operator of Monmouth Park Racetrack, the en banc Third Circuit, 9--3, held that the 2014 partial repeal ``authorized'' gambling and that PASPA did not commandeer because it required no affirmative state action. The Supreme Court granted certiorari; the case was recaptioned after Phil Murphy became governor.

\section{Judicial Votes}

Reversed: \textbf{7--2} on commandeering, \textbf{6--3} on severability (May 14, 2018).
\begin{itemize}
\item \textbf{Majority:} Justice Alito, joined by Chief Justice Roberts and Justices Kennedy, Thomas, Kagan, and Gorsuch; Justice Breyer joined all but Part VI--B.
\item \textbf{Concurring:} Justice Thomas.
\item \textbf{Concurring in part and dissenting in part:} Justice Breyer.
\item \textbf{Dissenting:} Justice Ginsburg, joined by Justice Sotomayor and in part by Justice Breyer.
\end{itemize}

\section{Holding}

Congress may not issue direct orders to state legislatures, whether compelling them to enact or enforce a federal program or prohibiting them from enacting, modifying, or repealing state laws. A valid preemption provision must exercise an enumerated power and be best read as regulating private actors. PASPA's bars on state ``authorization'' and ``licensing'' command state legislatures and are unconstitutional; because no remaining provision is severable, the whole statute is invalid. The remark that Congress may regulate sports gambling directly was dicta.

\section{Analysis of Opinions}

\textbf{Majority (Alito).} New Jersey's partial repeal ``authorized'' gambling, since repealing a ban both permits the activity and empowers operators. Under \textit{New York v.\ United States} and \textit{Printz}, Congress regulates individuals, not States; the rule protects liberty, political accountability, and state budgets. The line between affirmative commands and prohibitions was ``empty.'' \textit{South Carolina v.\ Baker} and \textit{Reno v.\ Condon} involved evenhanded regulation; \textit{Hodel} and \textit{FERC} involved choices. Preemption of every kind works through laws regulating private actors, which \S~3702(1) is not. Severing the rest would invert Congress's 1992 priorities: leaving \S~3702(2) alone would yield a ``perverse policy'' making private sports gambling federally unlawful only where a State legalizes it, and would overwhelm a civil-injunction scheme scored at ``no cost'' to the federal government.

\textbf{Thomas (concurring).} He doubted Congress may ban purely intrastate gambling and urged revisiting severability doctrine, which rests on hypothetical intent that never passed bicameralism and presentment.

\textbf{Breyer (in part).} Section 3702(2) validly regulates private interstate commerce and should survive, even if that made New Jersey's victory ``mostly Pyrrhic.''

\textbf{Ginsburg (dissenting).} Assuming \S~3702(1) invalid, the Court should have used a scalpel, not ``an ax''; the remaining bans were valid commerce regulation under \textit{Gonzales v.\ Raich}.

\section{Examples: Future Applications}

\begin{enumerate}
\item \textbf{Same-Side: Immigration Information.} A law forbidding States to bar police from sharing immigration status with federal agents dictates state policy without regulating private conduct.
\item \textbf{Same-Side: Cannabis Freeze.} A law forbidding States to repeal their cannabis penalties locks state legislatures in place. Congress may criminalize cannabis directly under the Controlled Substances Act, but freezing state codes blurs voter accountability.
\item \textbf{Opposite-Side: Federal Crime.} Making private sportsbooks a federal felony regulates private actors directly.
\item \textbf{Opposite-Side: Evenhanded Privacy Law.} A ban on selling biometric records that binds data brokers and state DMVs alike is valid under \textit{Reno v.\ Condon}.
\item \textbf{Fence-Sitter: Drone Rights.} A federal ``right to operate'' drones plus a ban on state zoning and privacy laws looks like express preemption---unless the FAA supplies virtually no substantive federal regime for private actors.
\end{enumerate}

\section{Critique}

\textbf{Closing a loophole.} Federalist scholars such as Ilya Somin praised the Court for rejecting the mandate--prohibition divide, without which Congress could gut anti-commandeering by forbidding repeals. Justice Thomas's concurrence exposed a textualist tension over severability by ``hypothetical congressional intent.''

\textbf{Doubts and ironies.} Steven Schwinn argued that the rule against commandeering lacks textual grounding and should yield to political safeguards, as in \textit{Garcia}; others echoed Ginsburg that striking all of PASPA was a policy choice. Yet progressive jurisdictions quickly used \textit{Murphy} to defend sanctuary and cannabis policies.

\textbf{Competing values.} The case balances state sovereignty, political accountability, and fiscal discipline---keeping Congress from conscripting state legislatures without bearing the political or financial cost---against national regulatory flexibility and cooperative efficiency, which let Congress set national policy without building a large federal enforcement apparatus.

\section{Key Quotations}

\begin{quote}
``The basic principle --- that Congress cannot issue direct orders to state legislatures --- applies in either event.''
\end{quote}
\noindent --- Justice Alito, for the Court.

\begin{quote}
``A more direct affront to state sovereignty is not easy to imagine.''
\end{quote}
\noindent --- Justice Alito, for the Court.

\begin{quote}
``Congress can regulate sports gambling directly, but if it elects not to do so, each State is free to act on its own.''
\end{quote}
\noindent --- Justice Alito, for the Court.

\end{document}
